\documentclass[aspectratio=169]{beamer}
\usetheme{Madrid}
\usecolortheme{default}
\usepackage[utf8]{inputenc}
\usepackage{amsmath, amssymb}
\usepackage{graphicx}
\usepackage{tcolorbox}

\definecolor{UCBblue}{RGB}{0, 51, 153}
\setbeamercolor{palette primary}{bg=UCBblue,fg=white}
\setbeamercolor{palette secondary}{bg=UCBblue,fg=white}
\setbeamercolor{palette tertiary}{bg=UCBblue,fg=white}
\setbeamercolor{titlelike}{bg=UCBblue,fg=white}
\setbeamercolor{item}{fg=UCBblue}

\title[Nivelación - Sesión 3]{Taller de Nivelación / Introducción a la Mecatrónica:\\ Álgebra Lineal Aplicada a la Teoría de Redes}
\subtitle{Módulo 4: Dinámica Frecuencial y Visualización Logarítmica}
\author[M.Sc. Ing. B. Quiroga Turdera]{M.Sc. Ing. Bernardo Quiroga Turdera}
\institute[UCB Tarija]{Carrera de Ingeniería Mecatrónica\\ Universidad Católica Boliviana "San Pablo"}
\date{30 de Julio de 2026}

\begin{document}

\begin{frame}
    \titlepage
\end{frame}

\begin{frame}{Agenda de la Sesión}
    \tableofcontents
\end{frame}

\section{Reactancias Dinámicas}
\begin{frame}{La Dependencia de la Frecuencia ($\omega$)}
    En la sesión anterior analizamos fasores a una frecuencia estática. Sin embargo, en mecatrónica las señales cambian (aceleración de motores, ruido de sensores).
    
    La reactancia ($X$) no es constante; depende de la frecuencia angular ($\omega = 2\pi f$):
    \vspace{0.3cm}
    \begin{itemize}
        \item \textbf{Bobina (Inductor):} $\mathbf{Z}_L = j\omega L$. A mayor frecuencia, \textit{mayor} oposición. Bloquea cambios rápidos.
        \item \textbf{Capacitor:} $\mathbf{Z}_C = -j\frac{1}{\omega C}$. A mayor frecuencia, \textit{menor} oposición. Deja pasar el ruido de alta frecuencia.
    \end{itemize}
\end{frame}

\section{El Problema del Espectro}
\begin{frame}{¿Por qué necesitamos computación modular?}
    Si queremos diseñar un filtro para nuestro robot, no basta con calcular la matriz $\mathbf{Z}$ a $50\text{ Hz}$. Debemos evaluar el sistema desde $1\text{ Hz}$ hasta $10,000\text{ Hz}$.
    
    \vspace{0.3cm}
    \begin{alertblock}{La Barrera del Cálculo Analítico}
    Resolver a pulso un sistema fasorial para $1,000$ frecuencias distintas es humanamente inviable. Es aquí donde la abstracción de la \textbf{Programación Modular (Funciones)} justifica su existencia en la ingeniería.
    \end{alertblock}
\end{frame}

\section{Visualización y Escalas Logarítmicas}
\begin{frame}{Representando el Espectro: Ejes Semilogarítmicos}
    Si graficamos de $1$ a $10,000\text{ rad/s}$ en un eje lineal estándar, las frecuencias bajas se comprimirán en el primer milímetro del gráfico, ocultando dinámicas vitales (como la resonancia).
    
    \vspace{0.3cm}
    \textbf{La solución (Gráficos de Bode preliminares):}
    Utilizamos escalas \textit{logarítmicas} para el eje X (Frecuencia).
    \begin{itemize}
        \item \textbf{Década:} Un intervalo donde el valor se multiplica por $10$ (de $1$ a $10$, de $10$ a $100$, etc.). Cada década ocupa el mismo espacio visual.
    \end{itemize}
\end{frame}

\section{Transición Computacional}
\begin{frame}{Herramientas en Python para el Análisis Espectral}
    \textbf{Fase CDIO: Diseñar e Implementar}
    
    En el \texttt{Notebook\_04}, integraremos tres herramientas clave:
    \begin{enumerate}
        \item \textbf{Modularidad (\texttt{def}):} Empaquetaremos nuestra matriz $\mathbf{Z}$ y el solver \texttt{np.linalg.solve} dentro de una función dependiente de $\omega$.
        \item \textbf{Vectores Logarítmicos (\texttt{np.logspace}):} Generaremos arreglos de prueba espaciados por décadas.
        \item \textbf{Dashboarding (\texttt{plt.semilogx}):} Utilizaremos \texttt{Matplotlib} para plotear la magnitud de la corriente y encontrar el punto de resonancia visualmente.
    \end{enumerate}
\end{frame}

\end{document}
